\section{Scope, thesis, and proof status}

\subsection{What the requested formalization should mean}

The phrase ``formalize algebraic geometry in iPRA'' can be understood in several inequivalent
ways.  The strongest and least useful interpretation would demand a single natural number coding
an arbitrary scheme, an internal set of all of its points, and a universal evaluator for every
algorithm mentioned in a textbook.  That interpretation is neither necessary nor compatible with
the foundational discipline of primitive recursive arithmetic.

The interpretation adopted here is presentation-relative and theorem-relative:

\begin{enumerate}[label=\arabic*.]
\item Fix a class of finite codes: polynomials, matrices, finite presentations, local fractions,
      overlap maps, and certificates.
\item Fix each primitive-recursive operation before it appears in a theorem.  There is no evaluator
      for arbitrary encoded programs.
\item Express equality and side conditions through finite witnesses whose checkers are primitive
      recursive.
\item Replace an infinite geometric object by the finite algebraic presentation needed by the
      theorem: an affine presentation, a finite atlas, a finite cover, a finite \v{C}ech complex, or
      finitely many observations of a represented object.
\item Prove the finite correctness statement with quantifier-free induction and the ordinary
      intuitionistic logical rules of iPRA.
\end{enumerate}

A completed formalization is not merely a Lean theorem about standard natural numbers.  In the
repository's notation, its endpoint is a term
\[
                 \Prf\ \PRAAx\ []\ \varphi,
\]
where $\varphi$ is the coded first-order formula expressing the advertised result.  The existing
bridge \code{pra_iff_native} may then transport that derivation to the concise native calculus.

\subsection{The central thesis}

The topological point-set definition of a scheme quantifies over prime ideals and arbitrary open
sets.  Most concrete arguments, however, consume much less information.  They use finitely
presented rings, standard opens $D(f)$, localizations, finitely many overlap maps, and finitely many
compatibility equations.  This finite algebraic skeleton is enough for a large body of algebraic
geometry.

The central thesis is therefore:

\begin{quote}
\emph{For iPRA, the right primitive notion is not an internally completed prime spectrum.  It is a
certified finite affine atlas together with proof-relevant localization and descent data.}
\end{quote}

The conventional scheme is the external mathematical interpretation of that atlas.  The internal
arithmetic proof manipulates only its finite codes and certificates.  This is analogous to working
with a finite presentation of a group rather than with a completed set containing all group
elements, but it additionally records the local-to-global structure that distinguishes geometry
from bare commutative algebra.

\subsection{Main contributions}

The report develops the following chain.

\begin{enumerate}[label=\textbf{C\arabic*.}]
\item \textbf{Certified setoids.}  Equality is a primitive-recursive relation
      $\EqW(x,y,w)$ on two object codes and a witness code.  Reflexivity, symmetry,
      transitivity, and operation congruence are implemented by witness transformers.  No equality
      search is assumed.
\item \textbf{Presented algebra.}  Equality in
      $\Z[X_0,\ldots,X_{n-1}]/(g_0,\ldots,g_{r-1})$ is witnessed by polynomials
      $h_i$ satisfying $p-q=\sum_i h_i g_i$.  Morphisms are polynomial tuples with relation
      certificates.
\item \textbf{Localization.}  Equality of $a/f^n$ and $b/f^m$ is witnessed by an exponent $k$
      and a base-ring certificate for
      $f^k(f^m a-f^n b)=0$.  The report gives explicit witness formulas for transitivity and ring
      congruence.
\item \textbf{Power covers.}  From $\sum_i c_i f_i=1$ and $N\geq1$, a bounded multinomial
      construction produces $d_i$ with $\sum_i d_i f_i^N=1$.
\item \textbf{Zariski gluing.}  Compatible local fractions in $M_{f_i}$ are homogenized, their
      overlap annihilators are synchronized, and one explicit linear combination produces a global
      section.  The same power-cover lemma proves uniqueness.
\item \textbf{Finite affine atlases.}  Charts, principal overlaps, localization isomorphisms, and
      cocycles form a finite presentation of a scheme.  Quasi-coherent modules and vector bundles
      become finite descent data.
\item \textbf{Projective geometry.}  The standard affine charts of $\Proj$ are explicit
      localizations.  For $\mathbf P^1$, Laurent-polynomial support gives complete normal forms for
      the \v{C}ech groups of $\cO(d)$.  For $\mathbf P^n$, exponent-vector summands admit an
      explicit finite contracting homotopy.
\end{enumerate}

\subsection{What is and is not proved}

The ordinary mathematical proofs of the power-cover lemma, localization calculus, principal-cover
gluing, and the projective \v{C}ech computations are complete in this report.  Their algorithms are
implemented over exact rings in the Python companion and checked on finite test suites, including a
zero-divisor example over $\Z/12\Z$.

The supplied Lean files are blueprints.  They contain host-level structures, definitions, and named
target propositions, together with the exact object-theory endpoint.  The preparation environment
had no Lean executable or local ProveIt dependency closure, so no new compilation claim is made.
The final conversion of the principal theorems to \code{Prf PRAAx} remains an engineering task,
although the required primitive recursions and quantifier-free invariants are specified.

The framework does not manufacture a general height bound, a complete rational-point list, a set of
Mordell--Weil quotient representatives, or an exact branch selector for arbitrary Cauchy
coefficients.  It makes such data composable and verifiable when supplied; it does not erase the
mathematics needed to produce them.
